% =====================================================================
% TAREA 09 - DOCUMENTACION DE LA CALIDAD (Semana 11, Unidad III)
% Proyecto Integrador: Outline (Wiki interna estilo Notion)
% Grupo TIID05B - Equipo 4
% ---------------------------------------------------------------------
% SECCION 4.6 "Aplicacion"  ->  Brian Odin Rubio Viramontes (Developer)
% Formato solicitado: Texto + Tabla de resultados. Citas: ISO 25010, IEEE 730.
% Basado en datos REALES del Examen Practico U2 (P2) del equipo.
% =====================================================================

\documentclass[11pt,letterpaper]{article}

\usepackage[utf8]{inputenc}
\usepackage[T1]{fontenc}
\usepackage[spanish]{babel}
\usepackage[margin=2.5cm]{geometry}
\usepackage{float}   % opcion [H] para fijar la tabla
\usepackage{array}   % columnas tipo p{}
\usepackage{hyperref}

\renewcommand{\tablename}{Tabla} % "Tabla" en lugar de "Cuadro"

\title{Tarea 09 --- Documentaci\'on de la Calidad de Outline\\[2pt]
       \large Secci\'on 4.6: Aplicaci\'on (borrador de Odin)}
\author{Brian Odin Rubio Viramontes --- TIID05B, Equipo 4}
\date{Julio 2026}

\begin{document}
\maketitle

% ---------------------------------------------------------------------
% NOTA PARA OLIVIA (unificacion en Overleaf):
%  - Este archivo contiene UNICAMENTE la seccion 4.6.
%  - Para integrarlo, copia el bloque entre los marcadores
%    "INICIO / FIN SECCION 4.6 (ODIN)".
%  - Uso \section*{} y \subsection*{} para que en el preview no salga
%    doble numeracion. Si quieres numeracion automatica, quita los
%    asteriscos (*) y ajusta los contadores segun el estilo del equipo.
% ---------------------------------------------------------------------

% =====================================================================
% ===================  INICIO SECCION 4.6 (ODIN)  =====================
% =====================================================================
\section*{4.6 Aplicaci\'on al proyecto Outline}

% ---------- INDICE COMENTADO ----------
\subsection*{\'Indice comentado del reporte}

A modo de esqueleto, el reporte de aplicaci\'on se organiza en los siguientes
apartados, cada uno con el contenido que le corresponde:

\begin{itemize}
  \item \textbf{Introducci\'on}: presenta el objetivo del reporte y describe
        brevemente qu\'e es Outline y por qu\'e se documenta su calidad.
  \item \textbf{Alcance}: delimita los m\'odulos de Outline sometidos a prueba
        y aquellos que quedaron fuera del an\'alisis.
  \item \textbf{Metodolog\'ia}: explica el enfoque empleado en el Examen
        Pr\'actico de la Unidad 2 (pruebas funcionales con Selenium y
        an\'alisis est\'atico con SonarQube).
  \item \textbf{Resultados}: consolida las m\'etricas cuantitativas obtenidas
        (cobertura, deuda t\'ecnica, casos de prueba y \textit{Quality Gate}).
  \item \textbf{Conclusiones}: sintetiza los hallazgos clave y las
        recomendaciones de remediaci\'on de la deuda t\'ecnica detectada.
\end{itemize}

% ---------- INTRODUCCION ----------
\subsection*{Introducci\'on}

El presente apartado documenta la aplicaci\'on pr\'actica de las estrategias de
aseguramiento de la calidad sobre el proyecto integrador \textbf{Outline}, una
plataforma \textit{open source} de gesti\'on del conocimiento empresarial de
arquitectura similar a Notion. A diferencia de las secciones anteriores ---de
car\'acter conceptual---, aqu\'i se reportan resultados concretos: la forma en
que se audit\'o la calidad del sistema mediante pruebas funcionales
automatizadas y an\'alisis est\'atico de c\'odigo, as\'i como las m\'etricas
obtenidas de dicho ejercicio.

Los resultados se enmarcan en el modelo de calidad de producto
\textbf{ISO/IEC 25010:2023} ---que define caracter\'isticas como la
mantenibilidad, la fiabilidad y la eficiencia de desempe\~no--- y en el
est\'andar \textbf{IEEE 730}, que estructura el plan de aseguramiento de la
calidad del software (SQA). El objetivo es evidenciar, de manera trazable y
verificable, el estado de calidad del m\'odulo evaluado de Outline.

% ---------- ALCANCE ----------
\subsection*{Alcance}

El an\'alisis pr\'actico se concentr\'o en los flujos funcionales cr\'iticos de
Outline accesibles desde su interfaz web, ejecutados sobre una instancia
desplegada localmente:

\begin{itemize}
  \item \textbf{Autenticaci\'on}: ingreso mediante correo electr\'onico y el
        flujo de inicio de sesi\'on.
  \item \textbf{B\'usqueda de documentos}: recuperaci\'on de art\'iculos
        existentes y manejo de resultados vac\'ios.
  \item \textbf{Creaci\'on y edici\'on documental}: validaci\'on del flujo de
        alta de documentos dentro del \textit{wiki}.
\end{itemize}

Estos m\'odulos se modelaron mediante el patr\'on Page Object Model (POM), con
clases dedicadas (\texttt{base\_page}, \texttt{login\_page},
\texttt{search\_page} y \texttt{document\_page}).

Quedaron \textbf{fuera del alcance} de las pruebas automatizadas: la
colaboraci\'on en tiempo real, las integraciones con servicios externos (por
ejemplo, autenticaci\'on real v\'ia Slack o \textit{webhooks}), la totalidad del
\textit{backend} y el renderizado recursivo del \'arbol de navegaci\'on de
colecciones (analizado de forma te\'orica en la Unidad 1, pero no cubierto por
la suite E2E).

% ---------- METODOLOGIA ----------
\subsection*{Metodolog\'ia}

El aseguramiento de la calidad se abord\'o con un enfoque dual y complementario:

\begin{itemize}
  \item \textbf{An\'alisis din\'amico (pruebas funcionales E2E)}: se implement\'o
        una suite de pruebas de extremo a extremo con \textbf{Selenium
        WebDriver} y \textbf{pytest}, bajo el patr\'on Page Object Model. La
        sincronizaci\'on se resolvi\'o exclusivamente con esperas expl\'icitas
        (\texttt{WebDriverWait}), evitando pausas fijas, y se incorporaron
        pruebas basadas en datos (\textit{data-driven}) mediante el decorador
        \texttt{@pytest.mark.parametrize}. La cobertura se midi\'o con
        \texttt{pytest-cov}, exportando los reportes \texttt{coverage.xml} y
        \texttt{report.html}.
  \item \textbf{An\'alisis est\'atico}: se audit\'o el c\'odigo fuente con
        \textbf{SonarQube v10.6}, evaluando sus cinco dimensiones (Seguridad,
        Confiabilidad, Mantenibilidad, Cobertura y Duplicaci\'on) y aplicando un
        \textit{Quality Gate} como criterio autom\'atico de aprobaci\'on.
\end{itemize}

El entorno de ejecuci\'on se levant\'o localmente mediante Docker (contenedores
de Redis y PostgreSQL) sobre el dominio \texttt{local.outline.dev:3000}, y el
flujo se automatiz\'o con un \textit{pipeline} de Integraci\'on Continua en
GitHub Actions.

\vspace{4pt}
\noindent\textit{Nota:} conforme a la naturaleza documental de esta tarea, esta
secci\'on \textbf{no ejecuta nuevas pruebas}; consolida y documenta los
resultados previamente obtenidos en el Examen Pr\'actico de la Unidad 2 (P2).

% ---------- RESULTADOS + TABLA ----------
\subsection*{Resultados: m\'etricas del Examen Pr\'actico U2}

La Tabla~\ref{tab:metricas-p2} sintetiza las m\'etricas cuantitativas obtenidas
durante el Examen Pr\'actico de la Unidad 2, junto con su interpretaci\'on.

\begin{table}[H]
\centering
\caption{M\'etricas de calidad obtenidas en el Examen Pr\'actico de la Unidad 2 (P2).}
\label{tab:metricas-p2}
\begin{tabular}{|p{3.1cm}|p{2.4cm}|p{8.4cm}|}
\hline
\textbf{M\'etrica} & \textbf{Valor} & \textbf{Interpretaci\'on} \\
\hline
Cobertura de c\'odigo & 62\% &
Porcentaje de l\'ineas del paquete de pruebas ejercitadas por la suite,
reportado por \texttt{pytest-cov}. Se ubica por debajo del umbral ideal
($>$80\%), lo que se\~nala rutas de c\'odigo a\'un sin verificar. Se vincula con
la \textbf{Mantenibilidad} de ISO/IEC 25010. \\
\hline
Deuda t\'ecnica & 23 d\'ias &
Esfuerzo estimado por SonarQube para remediar la totalidad de los
\textit{code smells} detectados; equivale a jornadas de refactorizaci\'on
pendientes. \\
\hline
Pruebas pasadas & 42/42 &
Todos los casos ejecutados (42 de 42) resultaron exitosos (\textit{passed}),
validando funcionalmente los flujos de autenticaci\'on, b\'usqueda y creaci\'on
documental; adicionalmente, 8 casos quedaron omitidos (\textit{skipped}) por
dependencias del entorno. \\
\hline
Quality Gate & Aprobado &
El proyecto satisfizo los umbrales del \textit{Quality Gate} definido en
SonarQube; la puerta de calidad no bloque\'o la integraci\'on del c\'odigo. \\
\hline
\end{tabular}
\end{table}

En conjunto, las m\'etricas indican un sistema funcionalmente estable ---la
totalidad de los casos ejecutados pasaron--- que satisface el \textit{Quality
Gate} configurado. No obstante, es importante se\~nalar que las dimensiones de
\textbf{Seguridad} y \textbf{Confiabilidad} obtuvieron una calificaci\'on
\textit{rating} D, lo que revela una deuda t\'ecnica considerable (23 d\'ias) y
la necesidad de priorizar refactorizaciones en los \textit{sprints} venideros
para aproximar dichas dimensiones al est\'andar ideal A.

% =====================================================================
% ====================  FIN SECCION 4.6 (ODIN)  =======================
% =====================================================================


% ---------------------------------------------------------------------
% ESPACIO RESERVADO PARA LAS DEMAS SECCIONES DEL EQUIPO
% (borrar / completar en el otro chat o al unificar en Overleaf)
% ---------------------------------------------------------------------
%   Portada, Indice, Introduccion general, Conclusion, Referencias -> Olivia
%   4.1 Estandares  y  4.2 Anatomia   -> Alan
%   4.3 Trazabilidad y 4.5 Ejemplos   -> Omar
%   4.4 Metricas y KPIs               -> Noe
%   4.6 Aplicacion                    -> Odin  (ESTE ARCHIVO)
% ---------------------------------------------------------------------

\end{document}
